\chapter[Componentes]{componentes}
  Neste capítulo, são apresentados todos os componentes presentes no sldc,
incluindo a etapa de configuração do ambiente de programação e
do ambiente para a utilização do kinect, assim podendo ser aplicada e analisada por outros interessados.

\section{Contextualização}
  O principal incentivo em relação a esta proposta é fazer um produto que seja de grande ajuda
para a recuperação de indivíduos que apresentam deficiências motoras e ou físicas, e que não
seja de forma penosa e ou dolorosa, podendo ser feito em casa e de forma mais lúdica.

  Com isso o sldc busca analisar o movimento de um índividuo com um movimento previamente
armazenado, dando \textit{feedbacks} em tempo real, para a correta execução. Além disso,
o sldc tem como objetivo a implementação de análise do movimento através da primeira
versão do kinect por ser mais facilmente encontrado e não gera grandes custos, porém
a precisão do kinect não é perfeita, sofrendo também de limitação por movimentos mais complexos.

\section{Ambiente}
  Nesta seção, será apresentado o passo a passo para configuração e integração das
tecnologias utilizadas durante este trabalho. A principal tecnologia utilizada é o Unity3d,
o qual possibilita a utilização da linguagem C\# para implementação dos algoritmos durante a realização do trabalho.
Também é utilizado o \textit{kinect for windows SDK V1.8 } que permite aos desenvolvedores
criar aplicativos que suportem o uso do kinect em computadores com windows.

 O SDK é necessário, pois o sistema operacional reconhece o sensor e disponibiliza a comunicação com o Unity3d
 , possibilitando a execução de códigos referente ao kinect. Além disso, o SDK envolve uma variedade
de classes e métodos com algoritmos já produzidos e disponibilizados para utilização
durante o desenvolvimento. Com a utilização da SDK, o desenvolvedor não se preocupa
com a implementação de algoritmos comuns, como a obtenção de dados dos
sensores.
A tecnologia Unity3d é multiplataforma, possibilitando sua utilização em vários Sistemas operacionais
porém o SDK é somente suportado no windows. Nas seções a seguir apresentadas, são descritos os passos para
instalação e configuração de ambas tecnologias no windows.

\subsection{Instalação do SDK}
  Primeiramente, para possibilitar a comunicação entre o PC e o kinect, deve-se
instalar o SDK 1.8, o qual pode ser encontrado \href{https://www.microsoft.com/en-us/download/details.aspx?id=40278}{aqui}.
  Lembrando que o kinect é feito para console, não é comum computadores darem suporte para o formato do
cabo, sendo assim para comunicar o PC e Kinect será preciso de uma fonte para Kinect Xbox 360.
  Depois do sdk instalado, conecte o Kinect na fonte e a fonte no PC. Para verificar se seu sitema operacional
está comunicando com o sensor alguns passos devem ser seguidos:

\begin{enumerate}
  \item Clicar em iniciar
  \item Ir em Meu computador
  \item Clicar com o botão direito
  \item Clicar em propriedades
  \item Na janela que foi aberta, ir em Gerenciador de dispositivos
  \item Verificar se na lista de dispositivos se encontra "Kinect for Windows" exibido como na imagem abaixo.
\end{enumerate}


\begin{figure}[!h]
\centering
\includegraphics [keepaspectratio=true,scale=0.60]{figuras/gerenciadorDispositivo.eps}

\caption{Gerenciador de Dispositivo}
\label{gerenciadorDispositivo}
\end{figure}

\subsection{Instalação do Unity 3d}
  Como já dito em \ref{Sec:Interface humano/computador}, o Unity 3d é um motor gráfico, criado para jogos em três dimensões e duas dimensões,
sendo a principal tecnologia usado para o desenvolvimento do sldc, para cálculos das juntas e para o \textit{feedback} ao usuário.
O Unity 3d pode ser encontrado \href{https://unity3d.com/pt/get-unity/download}{aqui}.

  Ao instalar o unity 3d, o instalador dará
a opção de instalar o visual studio mais recente, é aconselhável aceitar pois este é um editor de texto com várias funções para a edição
de arquivos c\#.

\subsection{Utilizndo o Unity 3d}
  Após a instalação de ambas as tecnologias, na primeira vez em que se abre o unity 3d é necessário fazer um cadastro para sua utilização. Feito
o cadastro:
\begin{enumerate}
  \item Na primeira janela do unity 3d é perguntado se deseja criar um novo projeto ou importar, clique em importar
  \item Procure o local onde se encontra a pasta do projeto e importe.
\end{enumerate}

  Pronto, agora será possível mexer no projeto pelo unity 3d. Para abrir o projeto no visual studio é somente clicar com o botão direito em cima
da cena desejada e clicar em "\textit{Open C\# Project}".
